\section{Computability and category}\label{sec:preliminaries}

We first encode ordinary real arithmetic by Turing degrees.
The category lemma at the end of the section will remove the
auxiliary real parameters introduced by the recovery construction.

We identify $2^\N$ with the sets of natural numbers, equipped with
the usual Cantor topology. Fix an effective enumeration
$(\Phi_e)_{e\in\N}$ of oracle Turing machines. For sets $A,B$, their
join $A\oplus B$ interleaves their characteristic functions.
Finite joins use a fixed effective coding.

The degree join $\mathbf a\vee\mathbf b$ is the least upper bound
of $\mathbf a,\mathbf b$: an oracle computes $A\oplus B$ exactly
when it computes both $A$ and $B$. Every order automorphism
therefore preserves finite joins. It also fixes the least degree
$\mathbf0$.

To distinguish ordinary real numbers from oracle sets, fix an
effective listing $(q_i)_{i\in\N}$ of $\Q$ and put
\[
C(x)=\{i:q_i<x\},\qquad d(x)=\degT{C(x)}
\quad(x\in\R).
\]
Thus all computability assertions about real numbers below refer
to these strict rational cuts.

\begin{lemma}\label{lem:cuts}
The cut map $C:\R\to2^\N$ is Borel and injective. If
$c_0,\ldots,c_k\in\Q$ and $x_1,\ldots,x_k\in\R$, with $k\ge1$, then
\[
d\left(c_0+\sum_{i=1}^k c_ix_i\right)
\leT \bigvee_{i=1}^k d(x_i).
\]
Every noncomputable degree equals $d(t)$ for some irrational
$t\in(0,1)$.
\end{lemma}

\begin{proof}
The $i$th coordinate of $C$ is the indicator of the open ray
$(q_i,\infty)$, so $C$ is Borel. Density of $\Q$ gives injectivity.

An oracle for $C(x)$ computes rational brackets $a<x\le b$ of
arbitrarily small width: search over rational pairs and check their
cut bits. The input cuts consequently compute rational
approximations, with prescribed error, to any fixed rational
linear combination $z$. If $z$ is irrational, comparison with
each rational query eventually separates and computes $C(z)$.
If $z$ is rational, its cut is computable. This proves the degree
inequality. It asserts existence of a reduction for each fixed
tuple, and does not require an algorithm deciding whether the
combination is rational.

Given a noncomputable set $B$, let
$t=\sum_{n\ge0}B(n)2^{-n-1}$. Then $0<t<1$ and $t$ is irrational:
the endpoint expansions and all binary expansions of rational
numbers are computable. The bits of $B$ give effective
approximations to $t$, hence compute $C(t)$ by irrationality.
Conversely, successive dyadic comparisons with the cut recover
the unique binary expansion of $t$. Thus $C(t)\equivT B$.
\end{proof}

Recall that a set is \emph{meager} if it is a countable union of
nowhere-dense sets, and \emph{comeager} if its complement is
meager. We use the Baire Category Theorem and the fact that Borel
sets in Polish spaces have the Baire property
\cite[Theorem~2.52 and Corollary~2.57]{MarkerDST}.

\begin{lemma}\label{lem:continuity}
If $X$ is Polish, $Y$ is second countable, and $f:X\to Y$ is
Borel, there is a comeager set $D\subseteq X$ such that $f|D$
is continuous.
\end{lemma}

\begin{proof}
Let $(U_n)$ be a countable basis for $Y$. By the Baire property,
write $f^{-1}(U_n)\mathbin{\triangle}V_n\subseteq M_n$, where
$V_n$ is open and $M_n$ is meager. On
$D=X\setminus\bigcup_nM_n$ one has
$(f|D)^{-1}(U_n)=D\cap V_n$, which is relatively open.
\end{proof}

Only continuity on the indicated restriction is asserted.
For a discrete-valued function, Lemma~\ref{lem:continuity}
makes its value constant on the intersection of that restriction
with a suitable neighborhood of each of its points.

Lemma~\ref{lem:avoidance} will remove auxiliary real parameters from
a degree bound that holds on a comeager set of pairs.
It is the relative category cone-avoidance principle attributed
to Sacks in \cite[Lemma~4.1]{KumarShelah2023}. We include the version for
rational-cut oracles that our proof requires.

\begin{lemma}[Category avoidance]\label{lem:avoidance}
Let $A,Y\subseteq\N$ with $Y\not\leT A$. Then
\[
\{(u,v)\in\R^2:Y\leT A\oplus C(u)\oplus C(v)\}
\]
is meager.
\end{lemma}

\begin{proof}
For a fixed index $e$, let $S_e$ be the set of pairs with both
coordinates irrational for which
$\Phi_e^{A\oplus C(u)\oplus C(v)}$ is the characteristic
function of $Y$. Suppose $S_e$ is dense in a nonempty open
rectangle $R$ with rational endpoints. We show that $Y\leT A$,
with $e$ and the endpoints of $R$ as finite program constants.

On input $n$, search for a finite halting computation of
$\Phi_e(n)$ whose answers about $A$ are correct and whose
finitely many proposed cut answers hold throughout a nonempty
open subrectangle of $R$. This is an effective search relative
to $A$. Dovetail simulations with finite time bounds and finite
tables of proposed answers to the two variable oracles.
For a cut $C(u)$, answer $1$ to a rational $q$ imposes $q<u$,
whereas answer $0$ imposes $u\le q$. Such a finite table is
compatible throughout a nonempty open interval inside the
corresponding side of $R$ exactly when the maximum of its lower
bounds is smaller than the minimum of its upper bounds,
after including the corresponding endpoints of $R$ among these bounds.
This is decidable rational arithmetic, and the same check
applies to $v$. Repeated queries must have consistent answers.

The search terminates. A pair in $S_e\cap R$ has a halting
computation on $n$. Since its coordinates are irrational, none
of the finitely many queried rationals equals the relevant
coordinate. Its finite transcript therefore holds on an open
neighborhood inside $R$, and appears in the search.
Every transcript accepted by the search is sound: its open
subrectangle intersects the dense set $S_e$, where the same
deterministic computation returns $Y(n)$. The search thus
computes $Y(n)$ with oracle $A$, a contradiction.

It follows that $S_e$ is not dense in any nonempty rational open
rectangle. Since these rectangles form a basis, $S_e$ is nowhere
dense. The pairs with a rational coordinate form a countable
union of nowhere-dense lines. Adding these pairs and taking
the countable union over $e$ completes the proof.
\end{proof}

The reduction index in Lemma~\ref{lem:avoidance} may depend on
$(u,v)$. Its proof takes the union over all indices, so no
uniformity across the success set is required.
